\section{从相似度分数到 Softmax Attention}
\label{intro:11-softmax-memory}

上一章已经用三个候选位置演示过：\textbf{先计算相关性分数，再将 Value 按权重组合}。
那一章着重解释“为什么需要”，这一章把计算写成\textbf{可以检查尺寸的矩阵公式}，并回答一条最常见的疑问：
\textbf{为什么一次查询写 $K\vect q$，多个查询却写 $QK^\top$？}

\subsection{Softmax 把相似度变成一组权重}

假设一个 Query 对两条 Key 的匹配分数分别为 $1$ 和 $0.6$。
指数函数 $e^{\beta s}$ 将分数转换成正数，
再把它们除以总和，便能得到总和为 $1$ 的权重。
$\beta$ 是缩放系数，它控制分数之间的相对差异，
不是必须手工调节的“记忆开关”。

先从两个已经给定的分数开始：
\[
s_A=1,\qquad s_B=0.6.
\]
指数变换得到 $e^{\beta}$ 和 $e^{0.6\beta}$。为了让两个权重的总和等于 $1$，将每个指数值除以指数值之和：
\begin{align}
p_A&=\frac{e^{\beta}}{e^{\beta}+e^{0.6\beta}},\\
p_B&=\frac{e^{0.6\beta}}{e^{\beta}+e^{0.6\beta}}.
\end{align}
这就是 \textbf{Softmax}。

\begin{table}[H]
\centering
\caption{调整 $\beta$ 时的 Softmax 权重（两个分数为 $1$ 和 $0.6$）}
\begin{tabular}{@{}rrr@{}}
\toprule
$\beta$ & $p_A$：较高分数 & $p_B$：较低分数\\
\midrule
$0$ & $50.0\%$ & $50.0\%$\\
$2$ & $69.0\%$ & $31.0\%$\\
$5$ & $88.1\%$ & $11.9\%$\\
$10$ & $98.2\%$ & $1.8\%$\\
$15$ & $99.8\%$ & $0.2\%$\\
\bottomrule
\end{tabular}
\end{table}

\begin{idea}
Softmax 本身不改变或创造 Value；它只把分数转成\textbf{从每条 Value 读取多少}的比例。在本例中，$\beta$ 增大让权重向较高分数倾斜。
\end{idea}

\subsection{给 Query、Key、Value 各自一个角色}

先规定一个\textbf{单次查询}：Query 是列向量 $\vect x$，
有 $M$ 个候选 Key/Value 对，每条 Key 含 $d_k$ 个数，每条 Value 含 $d_v$ 个数。
所以 $K\in\R^{M\times d_k}$、$V\in\R^{M\times d_v}$，
而查询 $\vect x\in\R^{d_k}$。第 $i$ 条 Key 与 Query 点积，
得到第 $i$ 条匹配分数；将 $M$ 个分数做 Softmax，
然后用所得权重组合对应 Value。

给第 $i$ 条候选的权重为：
\begin{equation}
p_i=\frac{e^{\beta\vect{k}_i^\top\vect{x}}}{\sum_j e^{\beta\vect{k}_j^\top\vect{x}}}.
\end{equation}
最终输出就是各个 Value 的加权平均：
\begin{equation}
\boxed{f(\vect{x})=\sum_i p_i\vect{v}_i
=\sum_i\vect{v}_i\frac{e^{\beta\vect{k}_i^\top\vect{x}}}{\sum_j e^{\beta\vect{k}_j^\top\vect{x}}}.}
\end{equation}
矩阵写法：
\begin{equation}
\boxed{f(\vect{x})=V^\top\softmax(\beta K\vect{x}).}\label{eq:column-attention}
\end{equation}
这是\textbf{按匹配分数组合 Value} 的单 Query 写法。此时还没有使用 Transformer 中的多 Query 矩阵。

\subsection{从单个列 Query 到标准 Attention：每一步都展开}

\subsection*{本节符号：先知道数字代表什么}

读接下来的公式时，先把下列符号与“数量”对应起来：

\begin{center}
\small
\renewcommand{\arraystretch}{1.3}
\begin{tabularx}{\linewidth}{@{}p{18mm}X p{37mm}@{}}
\toprule
符号 & 意思 & 数字例子\\
\midrule
$M$ & \textbf{Key--Value 记忆的条数} & 两条：$M=2$\\
$d_k$ & \textbf{每个 Key 有几个数字}，Query 维度与它相同 & $(1,0)$：$d_k=2$\\
$d_v$ & \textbf{每个 Value 有几个数字}，也是输出向量的维度 & $(2,1,0)$：$d_v=3$\\
$N$ & \textbf{同时处理多少个 Query}；批量查询时使用 & 三个：$N=3$\\
$\beta$ & 调整匹配分数的缩放系数 & 如 $\beta=1$\\
\bottomrule
\end{tabularx}
\end{center}

注意：\textbf{$M$ 是条数，$d_k$ 和 $d_v$ 是一条向量的长度。}下标 $k$ 与 $v$ 分别代表 Key 和 Value，$d$ 来自 \emph{dimension}（维度）。



刚才只有\textbf{一个列向量 Query}。标准 Transformer 一次处理多个 Token，各自发出 Query。看似不同的两条公式，实际上只需要统一行列约定、扩展成多个 Query，再选用标准缩放系数。

\textbf{第一步：把查询 $\vect x$ 改写为 $\vect q$。}在这个单次查询的简化例子中，取 $\vect q=\vect x$，因此
\[
\vect y=V^\top\softmax(\beta K\vect q).
\]
这里 $K$ 有 $M$ 行，每行 $d_k$ 个数字；$\vect q$ 有 $d_k$ 个数字；$V$ 有 $M$ 行，每行 $d_v$ 个数字。输出 $\vect y$ 最终有 $d_v$ 个数字。

\begin{example}[title={字母换成数字后，究竟在算什么？}]
设有\textbf{两个候选位置}（$M=2$），每条 Key 是\textbf{二维}（$d_k=2$），每条 Value 是\textbf{三维}（$d_v=3$）：
\[
K=\begin{bmatrix}1&0\\0&1\end{bmatrix},\quad
\vect q=\begin{bmatrix}1\\0\end{bmatrix},\quad
V=\begin{bmatrix}2&1&0\\0&3&4\end{bmatrix}.
\]
令 $\beta=1$，就得到
\[
K\vect q=\begin{bmatrix}1\\0\end{bmatrix},\qquad
\softmax(K\vect q)\approx
\begin{bmatrix}0.731\\0.269\end{bmatrix}.
\]
最后用匹配权重组合 Value：
\[
\vect y=V^\top\softmax(K\vect q)
\approx
\begin{bmatrix}2&0\\1&3\\0&4\end{bmatrix}
\begin{bmatrix}0.731\\0.269\end{bmatrix}
=\begin{bmatrix}1.462\\1.538\\1.076\end{bmatrix}.
\]
匹配分数有 \textbf{2 个}，是因为 $M=2$；最终输出有 \textbf{3 个数字}，是因为 $d_v=3$。两种数量不能混淆。
\end{example}

\textbf{第二步：统一改成行向量表示。}对矩阵转置，有
\[
(AB)^\top=B^\top A^\top,\qquad
(K\vect q)^\top=\vect q^\top K^\top.
\]
向量的一列分数和转置成一行的相同分数，Softmax 归一化值完全一样（只是排布方向变了）。于是对输出一起转置：
\begin{align}
\vect y^\top
&=\bigl[V^\top\softmax(\beta K\vect q)\bigr]^\top\nonumber\\
&=\bigl[\softmax(\beta K\vect q)\bigr]^\top V\nonumber\\
&=\boxed{\softmax(\beta\vect q^\top K^\top)V.}
\label{eq:single-row-attention}
\end{align}
\textbf{注意目前仍只有一个 Query！}左侧输出只是从列向量变为行向量。

\textbf{第三步：许多 Query 依次堆成行。}令
\[
Q=\begin{bmatrix}
\vect q_1^\top\\\vdots\\\vect q_N^\top
\end{bmatrix}\in\R^{N\times d_k}.
\]
把刚才一个 Query 的计算重复 $N$ 次，便可一次写成
\begin{equation}
\boxed{Y=\softmax(\beta QK^\top)V.}
\label{eq:multi-query-attention}
\end{equation}
相应维度如下：
\[
\underbrace{Q}_{N\times d_k}
\underbrace{K^\top}_{d_k\times M}
=\underbrace{S}_{N\times M},\qquad
\underbrace{\softmax(\beta S)}_{N\times M}
\underbrace{V}_{M\times d_v}
=\underbrace{Y}_{N\times d_v}.
\]
此时 Softmax 是\textbf{每一行单独归一化}：每个 Query 都为 $M$ 条 Key 分配权重。

\textbf{第四步：采用标准 Transformer 的缩放系数。}
Transformer 通常对未单位归一化的 Q、K 使用 $1/\sqrt{d_k}$，写成
\begin{equation}
\boxed{\operatorname{Attention}(Q,K,V)=
\softmax\!\left(\frac{QK^\top}{\sqrt{d_k}}\right)V.}
\label{eq:attention-transformer}
\end{equation}

\begin{warningbox}
\textbf{严格区分这几件事：}
从列向量变成行向量是\textbf{转置恒等式}；把一个 Query 变成多个 Query 是\textbf{批量计算}；把 $\beta$ 改为 $1/\sqrt{d_k}$ 是\textbf{缩放选择}，不是转置推导出来的结果。
原文的随机单位向量记忆模型也可能采用另一种 $\beta$。比较两种公式的数值结果前，必须先统一向量归一化与缩放约定。
\end{warningbox}

\begin{idea}[title={一句话读懂两种式子}]
列向量写法是\textbf{先匹配，再在左侧乘 $V^\top$}；行向量写法是\textbf{先匹配，再在右侧乘 $V$}。它们计算的是相同的加权 Value 组合，只是输出和输入各旋转了九十度。
\end{idea}

\begin{center}
\begin{tikzpicture}[>=Latex, node distance=5mm,
 block/.style={draw=linegray,rounded corners=3pt,fill=paper,align=center,minimum height=10mm,font=\small}]
\node[block,minimum width=30mm] (q) {Query $Q$};
\node[block,right=17mm of q,minimum width=30mm] (k) {Keys $K$};
\node[block,right=17mm of k,minimum width=30mm] (v) {Values $V$};
\node[block,below=13mm of k,minimum width=40mm] (dot) {相似度 $QK^\top$};
\node[block,below=of dot,minimum width=40mm] (soft) {Softmax 权重};
\node[block,below=of soft,minimum width=40mm] (out) {加权求和：输出};
\draw[->,blue,thick] (q.south) |- (dot.west);
\draw[->,blue,thick] (k.south)--(dot.north);
\draw[->,blue,thick] (dot)--(soft);
\draw[->,blue,thick] (soft)--(out);
\draw[->,orange,thick] (v.south) |- (out.east);
\end{tikzpicture}
\end{center}

\begin{warningbox}
\textbf{一个重要区别：}我们这里只使用\textbf{当前可读的候选 Key/Value}；真实 Self-Attention 的 Key/Value 是根据当前 Token 表示动态投影出来的。第六部分的联想记忆专题还会讨论\textbf{固定权重形式的 Key/Value}，它只是另一条理解路径，并不表示两种模型完全相同。
\end{warningbox}
